\section{Introduction}\label{sec:intro}

A numerical target written into law looks verifiable. It names a number, a population and a deadline, and it seems to ask only that someone compute the corresponding statistic and compare it with the threshold. In effect, such a target is a legal requirement to produce an official statistic, written before that statistic has been defined. What counts as a member of the population, as a request, as service and as the moment of observation depends on conventions that the law usually leaves open. Someone has to settle them after the law is passed, and that someone may be the agency whose performance the target is meant to assess. Until they are settled, the question of whether the target has been met has no unique answer, and a law that appeared verifiable is not.

The date of observation is the clearest example. When a target is set on the stock of a waiting list, the stock has to be counted on a particular day, and the answer moves with the day. Statistical systems that publish waiting lists usually fix that day by rule. Social housing registers kept by English local authorities are reported as of 31 March, Norwegian kindergarten waiting lists as of 15 December, and waiting times in the English National Health Service at the end of each reporting period \citep{mhclg_table600,ssb13632,nhs_rtt_mar26}. A fixed date makes counts comparable from one year to the next; it does not make any particular date neutral. The childcare waiting list analysed in this paper shows how much the choice can matter. Between 2014 and 2017, the stock counted on a single day fell by 53.2\% when that day was in December and by between 17.5\% and 20.5\% when it was in any other quarter, so that the percentage decline measured in December is 3.03 times that measured in September. Measures of flow, dated by the request rather than by the count, are much less sensitive to the ending month: for them, the ratio of the December change to the September change lies between 0.75 and 1.11. A law that sets a target on a stock without fixing the date leaves the verdict to whoever chooses the day. That is a matter of indicator design, and it arises whatever the service and whatever the country.

Much of this is familiar in outline. \citet{chun2005goal} describe evaluative goal ambiguity, the interpretive leeway that the goals of public agencies allow when progress towards them is judged. \citet{bevan2006measured} show that a regime of targets rests on the assumption that what is measured is what matters, and document the gaming that followed in the English health service. \citet{bird2005performance} set out what a performance monitoring protocol must specify for an indicator to be interpretable, from its definition to the reporting of its uncertainty. Work on multiverse and specification curve analysis recommends reporting results across the set of defensible analytic choices rather than for a single preferred one \citep{steegen2016increasing,simonsohn2020specification}. Together, these studies establish that public goals can be ambiguous, that indicators respond to the uses made of them and that conclusions can depend on choices taken before any number is reported. They seldom have the administrative microdata needed to enumerate the conventions behind a particular statutory target and to compute the verdict under each of them. In specification analysis, moreover, the alternatives usually concern estimation; here they arise earlier, when a legal objective is turned into a population, an event and a date. What microdata add is an arithmetic demonstration, in a case where the law, the data and an official verdict can all be examined, that the legal verdict depends on the convention.

Our case is S\~ao Paulo, where three conditions that rarely occur together are met. First, the ambiguity is in the statutory text. Goal~5 of the city's 2015 Municipal Education Plan requires childcare for at least 75\% of children aged 0 to 3 years and 11 months or for 100\% of registered demand, whichever is greater \citep{saopaulo2015pme}. The first requirement counts children whether or not their families have applied, the second rests on requests the administration can observe, and the clause that joins them admits more than one reading. Second, the inputs exist at the level of detail the conventions require. The municipality publishes its waiting list on regular dates and keeps centralized individual records of requests and of enrolments, and the School Census distinguishes enrolment by school network, so each convention can be computed rather than only argued about. Third, there is an official verdict to check: in 2023, the municipality's review of the plan declared Goal~5 achieved \citep{saopaulo2023conferencia}.

We make explicit the six conventions that the law leaves open: which waiting list counts as registered demand, the date on which it is observed, the source of enrolment for the demand requirement, the school networks counted for population coverage, the population series used as the denominator, and the rule that combines the two requirements. The inputs are the published waiting-list series, the municipality's enrolment publications, the School Census, live birth registrations, and demographic censuses and projections. The main grid fixes the December reference date, enrolment across all networks for the population requirement and the reading under which only the larger requirement binds, and crosses the remaining choices into 12 combinations. An extended grid in the appendix adds alternative dates, enrolment universes and readings of the clause, for 144 combinations. For each combination, we compute the verdict in every year in which its inputs exist. Each combination is a convention an agency could adopt with the data available. The grid shows how far the verdict depends on those conventions; it is not a distribution of legally valid readings, and we do not weight its cells.

Across our measurement grid, Goal~5 compliance is specification-dependent in 2019--2022, whereas no specification in that grid indicates compliance in any evaluable year in 2010--2018; before September 2019, only the published waiting-list series is available. Of the 12 main combinations, seven indicate compliance in at least one evaluable year and five never do, and none of the 648 evaluable combination-year observations for 2010--2018 in the extended grid indicates compliance. These statements hold within the documented grid. A mechanically rescaled population projection, which we exclude under an explicit provenance criterion, would indicate compliance in 2018. The 2023 review declared the goal achieved and reported that its threshold had been reached since 2015, without publishing its formula. No specification in the grid reproduces that account. The ratio of December municipal enrolment to the sum of that enrolment and the published waiting list reproduces six of the eight values the review reports, but a 75\% threshold on that ratio is not the statutory obligation to serve all registered demand, and it does not identify which requirement binds. Applied to the same ratio, the reference date alone moves the first year above 75\% from 2015 with December counts to 2018 with September counts.

What follows from these results is a list of what a verifiable indicator has to fix: the definitions of an eligible request and of service, the time within which a request must be met, the date of observation, the enrolment universe, the source and vintage of the population data, and the rule that combines the two requirements. Nothing on that list is specific to S\~ao Paulo or to childcare, since it applies to any target set on a waiting list or on the coverage of an eligible population. Brazil's new National Education Plan, enacted in April 2026, requires childcare to meet 100\% of expressed demand without defining the term, and the indicators that will measure it are still being specified \citep{brasil2026pne,mec2026portaria701}.

Section~\ref{sec:institutions} describes the legal and institutional setting, Section~\ref{sec:data} the data and Section~\ref{sec:goal} the verdict across the grid. Sections~\ref{sec:series} to~\ref{sec:service} examine the published stock, the coverage of the request record and the timing of requests and referrals, Section~\ref{sec:scope} sets out the limits of the evidence, Section~\ref{sec:indicators} turns the findings into indicators and a publication protocol, and Section~\ref{sec:conclusion} concludes.
